\chapter{Esercizi: Introduzione a Java}
\label{chap:esercizi-introduzione-java}

\begin{exercise}
\label{ex:fi-esercizi-introduzione-java-portabilita-cpu}
Perché Java viene progettato per essere indipendente dalla CPU?
\end{exercise}

Il sorgente viene compilato in bytecode destinato alla JVM, anziché direttamente
nelle istruzioni di una particolare CPU. Una JVM compatibile può eseguire lo
stesso bytecode su architetture differenti (\cref{sec:fi-java-esecuzione}).

\begin{exercise}
\label{ex:fi-esercizi-introduzione-java-compile-once}
Che cosa significa l'espressione ``compile once, execute everywhere''?
\end{exercise}

Il programma può essere compilato una volta in bytecode ed eseguito su
piattaforme diverse senza ricompilarlo. È comunque necessaria una JVM compatibile
e devono essere disponibili le risorse e le librerie richieste
(\cref{sec:fi-java-esecuzione}).

\begin{exercise}
\label{ex:fi-esercizi-introduzione-java-portabilita-confronto}
Perché Java è considerato più portabile di molti linguaggi compilati tradizionali?
\end{exercise}

Un eseguibile nativo contiene istruzioni per una specifica architettura; il
bytecode Java è invece destinato alla JVM, che ne gestisce l'esecuzione sulla
piattaforma ospitante (\cref{def:fi-java-jvm,sec:fi-linguaggi-esecuzione}).

\begin{exercise}
\label{ex:fi-esercizi-introduzione-java-platform-api}
Che cos'è la Java Platform API?
\end{exercise}

È l'interfaccia documentata delle classi e dei metodi forniti dalla piattaforma
Java: descrive le operazioni disponibili, i parametri, i risultati e, quando
pertinenti, le eccezioni
(\cref{sec:fi-java-invocazione-metodi,sec:fi-java-metodi-string}).

\begin{exercise}
\label{ex:fi-esercizi-introduzione-java-libreria-standard}
Perché la libreria standard di Java è importante per il programmatore?
\end{exercise}

Fornisce classi e metodi già disponibili per operazioni comuni, evitando di
doverle implementare nuovamente. Per utilizzarli occorre conoscere la relativa
interfaccia pubblica (\cref{sec:fi-java-classi-oggetti}).

\begin{exercise}
\label{ex:fi-esercizi-introduzione-java-sorgente-libreria}
Perché non è necessario conoscere il codice sorgente delle classi standard di Java?
\end{exercise}

È sufficiente conoscerne il contratto pubblico: metodi accessibili, parametri,
risultati e comportamento previsto. I dettagli interni di implementazione non
sono necessari all'invocazione dei metodi
(\cref{sec:fi-java-classi-oggetti,sec:fi-java-metodi-string}).

\begin{exercise}
\label{ex:fi-esercizi-introduzione-java-jdk}
Che cos'è il JDK?
\end{exercise}

Il \emph{Java Development Kit} è l'insieme di strumenti per sviluppare programmi
Java; comprende, fra gli altri, il compilatore \texttt{javac} e il comando
\texttt{java} per avviarli (\cref{sec:fi-java-esecuzione}).

\begin{exercise}
\label{ex:fi-esercizi-introduzione-java-fasi-programmazione}
Quali sono le tre fasi principali dell'attività di programmazione?
\end{exercise}

Nel ciclo di lavoro presentato negli appunti si distinguono \emph{modifica del
sorgente}, \emph{compilazione} e \emph{collaudo}. Eventuali errori richiedono di
ripetere il ciclo (\cref{fig:fi-java-ciclo-sviluppo,sec:fi-java-errori}).

\begin{exercise}
\label{ex:fi-esercizi-introduzione-java-codice-sorgente}
Che cosa si intende per codice sorgente?
\end{exercise}

È il testo del programma scritto dal programmatore in un linguaggio di
programmazione, prima della sua eventuale traduzione in una rappresentazione
eseguibile (\cref{sec:fi-linguaggi-esecuzione,sec:fi-java-esecuzione}).

\begin{exercise}
\label{ex:fi-esercizi-introduzione-java-estensione-java}
Che estensione deve avere un file sorgente Java?
\end{exercise}

L'estensione è \texttt{.java}, come in \texttt{HelloTester.java}
(\cref{sec:fi-java-esecuzione}).

\begin{exercise}
\label{ex:fi-esercizi-introduzione-java-javac}
Che cosa fa il comando \texttt{javac HelloTester.java}?
\end{exercise}

Invoca il compilatore Java sul file sorgente indicato; se la compilazione riesce,
produce il bytecode delle classi compilate (\cref{sec:fi-java-esecuzione}).

\begin{exercise}
\label{ex:fi-esercizi-introduzione-java-file-class}
Quale file viene generato compilando \texttt{HelloTester.java}?
\end{exercise}

La compilazione dell'esempio produce normalmente
\texttt{HelloTester.class}, contenente il bytecode della classe
(\cref{sec:fi-java-esecuzione}).

\begin{exercise}
\label{ex:fi-esercizi-introduzione-java-bytecode}
Che cos'è il bytecode Java?
\end{exercise}

È una rappresentazione intermedia del programma costituita da istruzioni
destinate alla JVM, non direttamente alla CPU fisica (\cref{def:fi-java-jvm}).

\begin{exercise}
\label{ex:fi-esercizi-introduzione-java-bytecode-cpu}
Perché il bytecode Java non è direttamente eseguibile dalla CPU reale?
\end{exercise}

Le sue istruzioni appartengono al modello di esecuzione della JVM, non
all'insieme di istruzioni macchina della CPU. Devono quindi essere interpretate o
tradotte in istruzioni native (\cref{def:fi-java-jvm}).

\begin{exercise}
\label{ex:fi-esercizi-introduzione-java-jvm}
Che cos'è la JVM?
\end{exercise}

La \emph{Java Virtual Machine} è l'ambiente che carica, verifica ed esegue il
bytecode Java, mediante interpretazione o compilazione in codice nativo
(\cref{def:fi-java-jvm}).

\begin{exercise}
\label{ex:fi-esercizi-introduzione-java-java-hello}
Che cosa fa il comando \texttt{java HelloTester}?
\end{exercise}

Avvia la classe \texttt{HelloTester} mediante il launcher Java, utilizzando la
JVM per eseguire il programma. Si indica il nome della classe, senza estensione
\texttt{.class} (\cref{sec:fi-java-esecuzione}).

\begin{exercise}
\label{ex:fi-esercizi-introduzione-java-interprete-ruolo}
Qual è il ruolo dell'interprete Java durante l'esecuzione di un programma?
\end{exercise}

L'interprete della JVM può eseguire le istruzioni del bytecode; la JVM può anche
compilarne alcune in codice nativo mediante JIT. L'interpretazione non è quindi
necessariamente l'unico meccanismo impiegato (\cref{sec:fi-java-esecuzione}).

\begin{exercise}
\label{ex:fi-esercizi-introduzione-java-compilato-interpretato}
Che differenza c'è tra un linguaggio compilato e un linguaggio interpretato?
\end{exercise}

Nella compilazione il codice viene tradotto prima dell'esecuzione in codice
macchina o intermedio; nell'interpretazione un ambiente esegue il sorgente o una
rappresentazione intermedia. I due meccanismi possono coesistere
(\cref{sec:fi-linguaggi-esecuzione}).

\begin{exercise}
\label{ex:fi-esercizi-introduzione-java-java-misto}
Perché Java può essere considerato un linguaggio ``misto''?
\end{exercise}

Il sorgente Java viene normalmente compilato in bytecode; durante l'esecuzione,
la JVM può interpretarlo oppure tradurlo in codice nativo tramite JIT. Si
combinano quindi compilazione ed esecuzione mediante macchina virtuale
(\cref{sec:fi-java-esecuzione}).

\begin{exercise}
\label{ex:fi-esercizi-introduzione-java-portabilita-sorgente-bytecode}
Che differenza c'è tra portabilità del codice sorgente e portabilità del bytecode?
\end{exercise}

La portabilità del sorgente riguarda la possibilità di compilare lo stesso
programma in ambienti diversi; quella del bytecode riguarda la possibilità di
eseguire lo stesso file compilato su JVM compatibili. Entrambe possono dipendere
dalle librerie e dall'ambiente richiesti (\cref{sec:fi-java-esecuzione}).

\begin{exercise}
\label{ex:fi-esercizi-introduzione-java-inizializzazione}
Che cosa significa inizializzare una variabile?
\end{exercise}

Significa assegnarle il primo valore, spesso contestualmente alla dichiarazione,
come in \texttt{int monete = 15000;}. Una variabile locale non inizializzata non
può essere letta prima di ricevere un valore (\cref{sec:fi-java-variabili}).

\begin{exercise}
\label{ex:fi-esercizi-introduzione-java-variabili-vantaggi}
Perché usare variabili rende un programma più leggibile e modificabile?
\end{exercise}

Un nome significativo identifica il ruolo di un dato e permette di riutilizzarne
il valore senza ripetere letterali o espressioni. Modificando un assegnamento si
possono aggiornare i calcoli che dipendono da quella variabile
(\cref{sec:fi-java-variabili,sec:fi-java-identificatori}).

\begin{exercise}
\label{ex:fi-esercizi-introduzione-java-identificatore-numero}
Perché un identificatore non può iniziare con un numero?
\end{exercise}

Le regole lessicali di Java non ammettono una cifra come primo carattere di un
identificatore, così da distinguere i nomi dai letterali numerici
(\cref{sec:fi-java-identificatori}).

\begin{exercise}
\label{ex:fi-esercizi-introduzione-java-parole-riservate}
Perché non si possono usare parole riservate come nomi di variabile?
\end{exercise}

Le parole riservate hanno un significato sintattico stabilito dal linguaggio, per
esempio \texttt{class} e \texttt{public}; non possono quindi essere riutilizzate
come identificatori (\cref{sec:fi-java-identificatori}).

\begin{exercise}
\label{ex:fi-esercizi-introduzione-java-nomi-simili}
Perché è sconsigliato usare nomi che differiscono solo per maiuscole e minuscole?
\end{exercise}

Java distingue maiuscole e minuscole, quindi i nomi sarebbero validi ma diversi.
La somiglianza aumenta il rischio di confonderli e rende meno leggibile il codice
(\cref{sec:fi-java-identificatori}).

\begin{exercise}
\label{ex:fi-esercizi-introduzione-java-separatori}
Quali sono alcuni separatori usati da Java?
\end{exercise}

Tra i simboli che delimitano costrutti o separano elementi figurano le parentesi
tonde \texttt{()}, graffe \texttt{\{\}}, quadre \texttt{[]}, il punto e virgola
\texttt{;}, la virgola e il punto. Per esempio, il punto separa l'oggetto dal
metodo invocato (\cref{sec:fi-java-sintassi,sec:fi-java-invocazione-metodi}).

\begin{exercise}
\label{ex:fi-esercizi-introduzione-java-operatori}
Quali sono alcuni operatori usati nelle espressioni Java?
\end{exercise}

Gli operatori aritmetici comprendono \texttt{+}, \texttt{-}, \texttt{*},
\texttt{/} e \texttt{\%}; \texttt{=} esegue un assegnamento e \texttt{==} un
confronto di uguaglianza. Altri operatori includono \texttt{<} e \texttt{>}
(\cref{sec:fi-java-variabili,sec:fi-java-tipi-numerici}).

\begin{exercise}
\label{ex:fi-esercizi-introduzione-java-file-nome-pubblicaclasse}
Perché il file deve chiamarsi esattamente \texttt{HelloTester.java}?
\end{exercise}

Una classe pubblica di primo livello chiamata \texttt{HelloTester}
deve trovarsi in \texttt{HelloTester.java}, rispettando le maiuscole e
le minuscole (\cref{sec:fi-java-primo-programma}).

\begin{exercise}
\label{ex:fi-esercizi-introduzione-java-case-sensitive}
Che cosa significa dire che Java è ``case-sensitive''?
\end{exercise}

Gli identificatori distinguono maiuscole e minuscole:
\texttt{HelloTester} e \texttt{hellotester} sono nomi diversi
(\cref{sec:fi-java-primo-programma}).

\begin{exercise}
\label{ex:fi-esercizi-introduzione-java-public-class-dichiarazione}
Che cosa definisce la riga \texttt{public class HelloTester}?
\end{exercise}

Dichiara una classe di nome \texttt{HelloTester} accessibile dall'esterno della
classe secondo le regole di visibilità del linguaggio
(\cref{sec:fi-java-primo-programma}).

\begin{exercise}
\label{ex:fi-esercizi-introduzione-java-public}
Che cosa indica la parola chiave \texttt{public}?
\end{exercise}

È un modificatore di accesso che rende una classe o un suo membro accessibile nei
contesti consentiti dal linguaggio; nella dichiarazione dell'esempio qualifica la
classe e il metodo \texttt{main} (\cref{sec:fi-java-primo-programma}).

\begin{exercise}
\label{ex:fi-esercizi-introduzione-java-class}
Che cosa indica la parola chiave \texttt{class}?
\end{exercise}

Introduce la dichiarazione di una classe, all'interno della quale possono essere
definiti campi e metodi
(\cref{sec:fi-java-primo-programma,sec:fi-java-classi-oggetti}).

\begin{exercise}
\label{ex:fi-esercizi-introduzione-java-main}
Che cos'è il metodo \texttt{main}?
\end{exercise}

Nella struttura tradizionale dei programmi Java adottata negli appunti, è il
metodo da cui comincia l'esecuzione dell'applicazione. La sua intestazione è
\texttt{public static void main(String[] args)}
(\cref{sec:fi-java-primo-programma}).

\begin{exercise}
\label{ex:fi-esercizi-introduzione-java-main-necessita}
Perché un'applicazione Java deve avere un metodo \texttt{main}?
\end{exercise}

Nella modalità di avvio tradizionale descritta negli appunti, la JVM deve
individuare un punto di ingresso da cui iniziare l'esecuzione; questo ruolo è
svolto dal metodo \texttt{main} con la firma prevista
(\cref{sec:fi-java-primo-programma}).

\begin{exercise}
\label{ex:fi-esercizi-introduzione-java-static-main}
Che cosa significa \texttt{static} nella dichiarazione del metodo \texttt{main}?
\end{exercise}

Indica che il metodo appartiene alla classe e può essere invocato senza creare
prima un oggetto della classe. Ciò permette di avviare il programma dal suo punto
di ingresso (\cref{sec:fi-java-primo-programma,sec:fi-java-invocazione-metodi}).

\begin{exercise}
\label{ex:fi-esercizi-introduzione-java-void-main}
Che cosa significa \texttt{void} nella dichiarazione del metodo \texttt{main}?
\end{exercise}

Indica che il metodo non restituisce alcun valore al chiamante
(\cref{sec:fi-java-primo-programma,sec:fi-java-metodi-string}).

\begin{exercise}
\label{ex:fi-esercizi-introduzione-java-args}
Che cosa rappresenta \texttt{String[] args}?
\end{exercise}

È il parametro del metodo \texttt{main} che può ricevere gli argomenti della riga
di comando come array di stringhe. Può rimanere inutilizzato nei programmi
semplici (\cref{sec:fi-java-primo-programma}).

\begin{exercise}
\label{ex:fi-esercizi-introduzione-java-commenti}
Che cosa sono i commenti in Java?
\end{exercise}

Sono parti del sorgente non eseguite come istruzioni: \texttt{//} introduce un
commento di una riga, mentre \lstinline|/*| e \lstinline|*/| delimitano un
commento potenzialmente su più righe (\cref{sec:fi-java-sintassi}).

\begin{exercise}
\label{ex:fi-esercizi-introduzione-java-commenti-utilita}
Perché i commenti sono utili anche se vengono ignorati dal compilatore?
\end{exercise}

Permettono di chiarire scelte progettuali, ipotesi e passaggi non immediati per
chi legge o modifica il programma. Non sostituiscono però nomi significativi e
codice comprensibile (\cref{sec:fi-java-sintassi,sec:fi-java-identificatori}).

\begin{exercise}
\label{ex:fi-esercizi-introduzione-java-system}
Che cos'è \texttt{System} in Java?
\end{exercise}

È una classe della libreria standard, \texttt{java.lang.System}, che fornisce
funzionalità collegate all'ambiente di esecuzione, tra cui l'accesso all'uscita
standard (\cref{sec:fi-java-uscita-standard}).

\begin{exercise}
\label{ex:fi-esercizi-introduzione-java-system-out}
Che cos'è \texttt{System.out}?
\end{exercise}

È un campo statico della classe \texttt{System} che contiene un riferimento a un
oggetto \texttt{PrintStream} associato all'uscita standard
(\cref{sec:fi-java-uscita-standard}).

\begin{exercise}
\label{ex:fi-esercizi-introduzione-java-println}
Che cosa fa il metodo \texttt{println()}?
\end{exercise}

Scrive il valore ricevuto sull'uscita standard e termina la riga. Una successiva
operazione di stampa inizia quindi dalla riga seguente
(\cref{sec:fi-java-uscita-standard}).

\begin{exercise}
\label{ex:fi-esercizi-introduzione-java-print-vs-println}
Qual è la differenza tra \texttt{print()} e \texttt{println()}?
\end{exercise}

Entrambi scrivono sull'uscita standard, ma \texttt{print()} non aggiunge un
terminatore di riga, mentre \texttt{println()} lo aggiunge
(\cref{sec:fi-java-uscita-standard}).

\begin{exercise}
\label{ex:fi-esercizi-introduzione-java-stringa}
Che cos'è una stringa in Java?
\end{exercise}

È una sequenza di caratteri rappresentata da un oggetto della classe
\texttt{String}. Un letterale si scrive tra virgolette doppie, per esempio
\texttt{"Hello, World!"}; gli oggetti \texttt{String} sono immutabili
(\cref{sec:fi-java-uscita-standard,sec:fi-java-metodi-string}).

\begin{exercise}
\label{ex:fi-esercizi-introduzione-java-invocare-metodo}
Che cosa significa ``invocare un metodo''?
\end{exercise}

Significa richiederne l'esecuzione, specificando il nome e gli eventuali
argomenti tra parentesi tonde. Il metodo esegue l'operazione prevista e può
restituire un risultato (\cref{sec:fi-java-invocazione-metodi}).

\begin{exercise}
\label{ex:fi-esercizi-introduzione-java-forma-invocazione}
Qual è la forma generale dell'invocazione di un metodo su un oggetto?
\end{exercise}

Per un metodo di istanza si utilizza \texttt{oggetto.nomeMetodo(argomento1,
argomento2)}, dove l'oggetto precede il punto e riceve l'invocazione
(\cref{sec:fi-java-invocazione-metodi}).

\begin{exercise}
\label{ex:fi-esercizi-introduzione-java-argomenti-separazione}
Come si separano più parametri in una chiamata di metodo?
\end{exercise}

Nella chiamata si forniscono gli \emph{argomenti}, separandoli con virgole
all'interno delle parentesi tonde, per esempio \texttt{Math.max(3, 7)}. I
\emph{parametri} sono invece dichiarati nell'intestazione del metodo
(\cref{sec:fi-java-invocazione-metodi,sec:fi-java-metodi-string}).

\begin{exercise}
\label{ex:fi-esercizi-introduzione-java-formato-libero}
Perché, anche se Java ha formato libero, è importante scrivere codice leggibile?
\end{exercise}

L'indentazione e la disposizione degli spazi non determinano i blocchi, ma
rendono riconoscibili la struttura e il flusso del programma. Il formato
leggibile facilita revisione e manutenzione (\cref{sec:fi-java-sintassi}).

\begin{exercise}
\label{ex:fi-esercizi-introduzione-java-errore-sintassi-logico}
Qual è la differenza tra errore di sintassi ed errore logico?
\end{exercise}

Un errore sintattico viola le regole del linguaggio e viene rilevato durante la
compilazione. Un errore logico può invece lasciare il programma eseguibile, ma
produrre risultati diversi da quelli richiesti (\cref{sec:fi-java-errori}).

\begin{exercise}
\label{ex:fi-esercizi-introduzione-java-system-aut}
Perché \texttt{System.aut.println("Hello, World!");} produce un errore di compilazione?
\end{exercise}

La classe \texttt{System} non contiene un campo chiamato \texttt{aut}: il nome
corretto è \texttt{out}. Il compilatore non riesce quindi a risolvere il simbolo
richiesto (\cref{sec:fi-java-errori}).

\begin{exercise}
\label{ex:fi-esercizi-introduzione-java-cannot-find-symbol}
Che cosa significa ``cannot find symbol'' in un messaggio del compilatore?
\end{exercise}

Indica che un nome utilizzato nel sorgente, per esempio una variabile, un campo o
un metodo, non è risolvibile nel contesto in cui compare. Non implica
necessariamente un errore di sola punteggiatura sintattica
(\cref{sec:fi-java-errori}).

\begin{exercise}
\label{ex:fi-esercizi-introduzione-java-errore-logico-runtime}
Che cosa distingue un errore logico da un errore di runtime?
\end{exercise}

Un errore logico produce un comportamento non conforme alla specifica, anche
senza interrompere il programma. Un errore di runtime si manifesta durante
l'esecuzione, per esempio come eccezione non gestita (\cref{sec:fi-java-errori}).

\begin{exercise}
\label{ex:fi-esercizi-introduzione-java-errore-file}
Fai un esempio di errore di runtime legato alla lettura di un file.
\end{exercise}

Un programma tenta di aprire un file inesistente: l'operazione può generare
un'eccezione, come \texttt{FileNotFoundException}, se non viene gestita. Il
sorgente può essere comunque sintatticamente corretto
(\cref{sec:fi-java-errori}).

\begin{exercise}
\label{ex:fi-esercizi-introduzione-java-errore-input}
Fai un esempio di errore di runtime legato all'input dell'utente.
\end{exercise}

Un programma che si aspetta un intero riceve testo non numerico e tenta di
convertirlo: può verificarsi un'eccezione, come \texttt{NumberFormatException}
nell'uso di \texttt{Integer.parseInt}
(\cref{sec:fi-java-errori,sec:fi-java-parsing}).

\begin{exercise}
\label{ex:fi-esercizi-introduzione-java-variabile}
Che cos'è una variabile in Java?
\end{exercise}

È un identificatore associato a un valore di un tipo dichiarato. Per un tipo
primitivo la variabile contiene il valore; per un tipo riferimento contiene un
riferimento a un oggetto (\cref{sec:fi-java-variabili,sec:fi-java-riferimenti}).

\begin{exercise}
\label{ex:fi-esercizi-introduzione-java-dichiarazione}
Che cosa significa dichiarare una variabile?
\end{exercise}

Significa indicarne il tipo e il nome, per esempio \texttt{int lit;}. La
dichiarazione può includere il valore iniziale, come in \texttt{int monete =
15000;} (\cref{sec:fi-java-variabili}).
